\documentclass[11pt]{article}

\usepackage[T1]{fontenc}
\usepackage{amsmath,amssymb,amsthm}
\usepackage[margin=1in]{geometry}
\usepackage[hidelinks]{hyperref}

\newtheorem{theorem}{Theorem}
\newtheorem{proposition}[theorem]{Proposition}
\newtheorem{lemma}[theorem]{Lemma}
\newtheorem{corollary}[theorem]{Corollary}
\newtheorem{conjecture}[theorem]{Conjecture}
\theoremstyle{remark}
\newtheorem{remark}[theorem]{Remark}

\newcommand{\R}{\mathbb R}
\newcommand{\T}{\mathbb T}
\newcommand{\Diss}{\mathcal{D}}
\newcommand{\WN}{\mathcal W^{\mathrm N}}

\title{Local NSE Sharpness of the Type-II Half-Radius Charge}
\author{John Robert Lawson and Codex}
\date{Research round ending 24 July 2026}

\begin{document}
\maketitle

\begin{abstract}
The preceding frozen-band theorem bounded all nonlinear correlation-death
events by one dissipation budget, but with a vanishing square-root-radius
charge.  This report proves that the half power cannot be improved by a
one-event estimate.  An explicit six-mode divergence-free field has
Gaussian-band energy \(651/2048\) and normalized nonlinear band flux
\(-109/2048\).  A solenoidal cutoff gives compact data on \(\R^3\).
Concentrating that datum at fixed energy and evolving by physical
Navier--Stokes produces, for every small radius \(r\), a smooth interval of
length \(r^{5/2}\) with fixed nonlinear frozen-band work and only
\(O(\sqrt r)\) dissipation.  The construction changes trajectory with
\(r\); it is a local sharpness family, not a singular cascade.
\end{abstract}

\section{Epistemic scope}

This round tests one possible closing move for the conditional Type-II
route: replacing the summable half-radius event charge by a stronger local
NSE estimate.  The result is negative.  The counterfamily consists entirely
of smooth compactly supported finite-energy data, but uses different initial
data at each radius.  It proves neither that one trajectory can concatenate
the packets nor that a first-record singular genealogy exists.

\section{Exact periodic seed}

On \(\T^3=(\R/2\pi\mathbb Z)^3\), let
\[
 \Phi(x,y)
 =
 \cos x+\cos(2y)+\cos(x+2y)
\]
and
\[
 U
 =
 (\partial_y\Phi,-\partial_x\Phi,0).
\]
Thus
\[
 U
 =
 \bigl(
 -2\sin(2y)-2\sin(x+2y),
 \ \sin x+\sin(x+2y),
 \ 0
 \bigr).
\]
The field is real, smooth, mean zero, and divergence-free.

For a radial Fourier multiplier \(M\), write
\[
 \mu_1=M(1),\qquad \mu_4=M(4),\qquad \mu_5=M(5),
\]
where the argument is squared frequency.  Let
\[
 \langle f\rangle_{\T^3}
 :=
 \frac1{|\T^3|}\int_{\T^3}f\,dx.
\]

\begin{proposition}[Exact triad flux]
\[
\boxed{
 \left\langle
 U\otimes U:\nabla(MU)
 \right\rangle_{\T^3}
 =
 -\frac12\mu_1+2\mu_4-\frac32\mu_5.
}
\]
\end{proposition}

\begin{proof}
The stream function has only the modes
\[
 \pm(1,0,0),\qquad
 \pm(0,2,0),\qquad
 \pm(1,2,0).
\]
For each mode,
\[
 \widehat U(k)
 =
 i(k_2,-k_1,0)\widehat\Phi(k).
\]
The normalized cubic mean is the finite convolution
\[
 \sum_{k+p+q=0}
 \sum_{\alpha,\beta=1}^3
 \widehat U_\alpha(k)
 \widehat U_\beta(p)
 (iq_\beta)\mu_{|q|^2}
 \widehat U_\alpha(q).
\]
Direct enumeration gives coefficients
\(-1/2,2,-3/2\) at squared frequencies \(1,4,5\).
The repository certificate performs this enumeration with exact rational
complex arithmetic rather than evaluating the closed formula.
\end{proof}

Define
\[
 S_*
 :=
 G_{\sqrt{\log2}}^2-G_{\sqrt{2\log2}}^2.
\]
Since \(G_r^2\) has multiplier \(e^{-r^2|\xi|^2}\), the multiplier at an
integer squared frequency \(n\) is
\[
 \mu_n=2^{-n}-4^{-n}.
\]
Hence
\[
 \mu_1=\frac14,\qquad
 \mu_4=\frac{15}{256},\qquad
 \mu_5=\frac{31}{1024}.
\]
The proposition yields
\[
\boxed{
 Q_{\T}
 :=
 \left\langle
 U\otimes U:\nabla(S_*U)
 \right\rangle_{\T^3}
 =
 -\frac{109}{2048}.
}
\]
Exact Parseval summation also gives
\[
\boxed{
 H_{\T}
 :=
 \left\langle U\cdot S_*U\right\rangle_{\T^3}
 =
 \frac{651}{2048}>0.
}
\]
Replacing \(U\) by \(-U\) preserves \(H_{\T}\) and reverses \(Q_{\T}\).
Thus the local quadratic NSE algebra has no universal sign obstruction to
Gaussian-band transfer.

\section{Compact localization}

Let
\[
 \mathcal A=(0,0,\Phi),
 \qquad
 \nabla\times\mathcal A=U.
\]
Choose \(\chi_L\in C_c^\infty(\R^3)\) equal to one on \([-L,L]^3\),
supported in its unit enlargement, with derivatives bounded independently
of \(L\).  Put
\[
 v_L:=\nabla\times(\chi_L\mathcal A).
\]
Then \(v_L\in C^\infty_{c,\sigma}(\R^3)\) and agrees with \(U\) in the
interior cube.

\begin{proposition}[Solenoidal localization preserves flux]
As \(L\to\infty\),
\[
 \frac1{|[-L,L]^3|}
 \int_{\R^3}
 v_L\otimes v_L:\nabla(S_*v_L)\,dx
 \longrightarrow Q_{\T},
\]
and
\[
 \frac1{|[-L,L]^3|}
 \langle v_L,S_*v_L\rangle
 \longrightarrow H_{\T}.
\]
\end{proposition}

\begin{proof}
Both \(S_*\) and \(\nabla S_*\) have Schwartz convolution kernels.
Choose an interior whose distance \(R_L\) from the cutoff tends to infinity
while \(R_L=o(L)\).  On that interior, the truncated convolution differs
from the periodic heat convolution by \(o(1)\), uniformly.  Its volume is
\(|[-L,L]^3|+o(L^3)\), so it contributes the asserted periodic mean.

The remaining layer has volume \(O(L^2R_L)=o(L^3)\), and the fields and
fixed derivatives are uniformly bounded there.  The kernel contribution
from beyond distance \(R_L\) is \(o(L^3)\) by Schwartz decay.  This proves
both limits.
\end{proof}

Fix a sufficiently large \(L\), reverse sign if desired, and write
\[
 v=v_L,\qquad
 H_*:=\frac12\langle v,S_*v\rangle>0,\qquad
 Q_*:=\int_{\R^3}v\otimes v:\nabla(S_*v)\,dx\ne0.
\]

\section{Fixed-energy concentration}

Let
\[
 \rho_{\mathrm L}=\sqrt{\log2},
 \qquad
 \rho_{\mathrm H}=\sqrt{2\log2},
\]
and
\[
 S_r=G_{\rho_{\mathrm L}r}^2-G_{\rho_{\mathrm H}r}^2.
\]
Given a desired band energy \(\gamma>0\), choose
\[
 A=\sqrt{\frac{\gamma}{H_*}}
\]
and define
\[
 u_r^0(x)=Ar^{-3/2}v(x/r).
\]
Gaussian covariance gives
\[
 S_ru_r^0(x)=Ar^{-3/2}(S_*v)(x/r).
\]
Therefore
\[
\boxed{
 \frac12\langle u_r^0,S_ru_r^0\rangle=\gamma,
}
\]
while
\[
 \frac12\|u_r^0\|_2^2
 =
 \frac{A^2}{2}\|v\|_2^2
\]
is independent of \(r\).  The critical endpoint size is
\[
\boxed{
 \|u_r^0\|_{L^{3,\infty}}
 =
 Ar^{-1/2}\|v\|_{L^{3,\infty}}.
}
\]

\section{Smooth NSE packet evolution}

Let \(u_r\) solve physical NSE with viscosity \(\nu>0\) and datum \(u_r^0\).
Set
\[
 y=x/r,\qquad
 s=\frac{A}{r^{5/2}}t,\qquad
 u_r(x,t)=Ar^{-3/2}V_r(y,s).
\]
Then
\[
 \partial_sV_r
 +\mathbb P\nabla_y\!\cdot(V_r\otimes V_r)
 =
 \varepsilon_r\Delta_yV_r,
 \qquad
 V_r(0)=v,
\]
with
\[
\boxed{
 \varepsilon_r=\frac{\nu\sqrt r}{A}\longrightarrow0.
}
\]

\begin{lemma}[Uniform short-time inviscid limit]
There is \(s_0>0\) such that the Euler solution \(V\) from \(v\) and all
sufficiently small-viscosity \(V_r\) exist smoothly on \([0,s_0]\), with
\[
 \sup_{0\le s\le s_0}
 \|V_r(s)-V(s)\|_{H^{m-1}}
 \lesssim\varepsilon_r
\]
for a fixed sufficiently large \(m\).
\end{lemma}

\begin{proof}
Smooth Euler and NSE energy estimates give a common \(H^{m+1}\) ceiling on
a short interval depending only on \(v\), uniformly for
\(0\le\varepsilon_r\le1\).  For \(w_r=V_r-V\), the difference equation and
standard commutator bounds give
\[
 \frac{d}{ds}\|w_r\|_{H^{m-1}}^2
 \le
 C\|w_r\|_{H^{m-1}}^2+C\varepsilon_r^2.
\]
Since \(w_r(0)=0\), Gronwall proves the estimate.
\end{proof}

The Euler pairing
\[
 q(s)
 :=
 \int_{\R^3}
 V(s)\otimes V(s):\nabla(S_*v)\,dy
\]
is continuous and \(q(0)=Q_*\ne0\).  Shrink \(s_0\) so that
\[
 |q(s)|\ge\frac{|Q_*|}{2}
 \quad(0\le s\le s_0)
\]
with constant sign.  The inviscid convergence then implies
\[
 \left|
 \int_0^{s_0}\!\!\int_{\R^3}
 V_r\otimes V_r:\nabla(S_*v)\,dy\,ds
 \right|
 \ge
 \frac{s_0|Q_*|}{4}
 =:w_*>0
\]
for small \(r\).

Set
\[
 t_r=\frac{s_0r^{5/2}}{A}.
\]
The physical nonlinear work against the frozen initial test is
\[
\begin{aligned}
 \WN_r
 &:=
 \int_0^{t_r}\!\!\int_{\R^3}
 u_r\otimes u_r:\nabla(S_ru_r^0)\,dx\,dt\\
 &=
 A^2
 \int_0^{s_0}\!\!\int_{\R^3}
 V_r\otimes V_r:\nabla(S_*v)\,dy\,ds.
\end{aligned}
\]
Consequently
\[
\boxed{
 |\WN_r|
 \ge
 A^2w_*
 =
 \frac{w_*}{H_*}\gamma
 =:\beta\gamma,
 \qquad \beta>0.
}
\]

\section{Half-radius dissipation}

On the same interval,
\[
\begin{aligned}
 \Diss_r
 &:=
 \nu\int_0^{t_r}\|\nabla u_r(t)\|_2^2\,dt\\
 &=
 \nu A\sqrt r
 \int_0^{s_0}\|\nabla V_r(s)\|_2^2\,ds.
\end{aligned}
\]
The carrier integral converges to a finite positive Euler integral.
Therefore
\[
\boxed{
 c\nu A\sqrt r
 \le
 \Diss_r
 \le
 C\nu A\sqrt r.
}
\]
The lower physical Gaussian radius is
\[
 r_r^{\mathrm L}=\rho_{\mathrm L}r,
\]
so this is exactly the half-radius scale of the common frozen-band budget.

The full monomial ledger is
\[
\begin{array}{c|c}
\text{quantity}&\text{power of }r\\ \hline
\text{velocity amplitude}&-3/2\\
\text{weak-}L^3\text{ size}&-1/2\\
\text{turnover time}&5/2\\
\text{nonlinear work rate}&-5/2\\
\text{enstrophy}&-2\\
\text{integrated nonlinear work}&0\\
\text{dissipation}&1/2\\
\text{effective viscosity}&1/2
\end{array}
\]
and uniform carrier bounds also give
\[
 \int_0^{t_r}
 \|u_r(t)\|_{L^{3,\infty}}^4\,dt
 \lesssim\sqrt r.
\]

\begin{theorem}[Local NSE sharpness of the half power]
For fixed \(\gamma,\nu>0\), there is a family of smooth compactly supported
finite-energy NSE data \(u_r^0\) with energy uniformly bounded as
\(r\downarrow0\), and smooth intervals \([0,t_r]\), such that:
\begin{enumerate}
\item the initial positive Gaussian band has energy \(\gamma\);
\item the nonlinear frozen-band work is at least \(\beta\gamma\);
\item \(t_r\asymp r^{5/2}\);
\item \(\Diss_r\asymp\nu\sqrt r\); and
\item the lower Gaussian radius is comparable to \(r\).
\end{enumerate}
\end{theorem}

\begin{corollary}[No stronger one-event radius charge]
For every \(\delta>0\), no estimate
\[
 \Diss([0,t])
 \ge
 c_{\gamma,\nu,E}(r^{\mathrm L})^{1/2-\delta}
\]
can hold uniformly for all smooth fixed-fraction band-work events using only
\(\gamma,\nu\), a kinetic-energy ceiling, and local NSE evolution.
\end{corollary}

\begin{proof}
The theorem would imply
\[
 c_{\gamma,\nu,E}r^{1/2-\delta}
 \lesssim\nu\sqrt r.
\]
After division by \(r^{1/2-\delta}\), the right side tends to zero like
\(r^\delta\).
\end{proof}

\section{Findings subject to review}

\begin{enumerate}
\item The exact Gaussian multipliers, band pairing, and nonlinear triad flux
are finite rational calculations and are executable.
\item Solenoidal localization preserves the nonzero periodic flux.
\item Physical NSE packets realise fixed nonlinear band work with
square-root dissipation and occupation.
\item The half-radius loss is locally sharp even for genuine smooth NSE,
not merely for an abstract Hilbert rotation.
\item Either local transfer sign occurs.
\end{enumerate}

The finite certificate does not verify the localization or local PDE
argument.  Those proofs are written above and await external mathematical
review.  No literature-novelty claim is made.

\section{Conjectures and remaining obligations}

\begin{conjecture}[Same-trajectory recycling obstruction]
One finite-energy first-record trajectory cannot concatenate the sharp local
packets while repeatedly reusing one fixed band-energy quantum without
paying a nonsummable transition, spatial-ancestry, or temporal-variation
charge.
\end{conjecture}

\begin{conjecture}[Unthinned genealogical Carleson law]
The heat-scale intervals and birth--death times of the complete unthinned
record family define a Carleson measure only after a same-trajectory
transition term is included; that term counts reuse hidden from the
one-event half-radius budget.
\end{conjecture}

The immediate obligations are:
\begin{enumerate}
\item compare the terminal carrier state of one sharp packet with the
initial carrier state required at the next smaller first record;
\item quantify the cost of spatially moving or spectrally reorienting that
state;
\item distinguish genuine reuse from independent varying-trajectory
packets;
\item keep coherent-trace rigidity and divergent-normalised-energy control
as separate branches.
\end{enumerate}

\section{Conclusion}

The local triadic escape is real and saturates the previous theorem.  The
remaining obstruction is therefore genuinely genealogical: it concerns
concatenating events on one trajectory, not strengthening a single-event
energy estimate.  This report proves no regularity or breakdown theorem and
no Clay alternative A--D.

The proof-owning note is stored as
\begin{center}
\small
\texttt{dossier/experiments/}\\
\texttt{type-ii-triad-packet-sharpness.md}.
\end{center}
The note was retired in the 2026-07-27 curation; it is recoverable from
Git history before commit \texttt{13557eb}, per the recipe in
\texttt{dossier/records/README.md}.
The exact certificate is
\texttt{lab/navier\_lab/type\_ii\_triad\_packet.py}.
The public repository is
\url{https://github.com/Bingham-Research-Center/brc-navier-stokes}.

\end{document}
